\section{AHSME 1994}
        \begin{enumerate}
        	\item (\textit{AHSME 1994}) $4^4 \cdot 9^4 \cdot 4^9 \cdot 9^9=$
$\textbf{(A)}\ 13^{13} \qquad\textbf{(B)}\ 13^{36} \qquad\textbf{(C)}\ 36^{13} \qquad\textbf{(D)}\ 36^{36} \qquad\textbf{(E)}\ 1296^{26}$

 \textbf{Answer: }C



\textbf{Solution}

Note that $a^x\times a^y=a^{x+y}$. So $4^4\cdot 4^9=4^{13}$ and $9^4\cdot 9^9=9^{13}$. Therefore, $4^{13}\cdot 9^{13}=(4\cdot 9)^{13}=\boxed{\textbf{(C)}\ 36^{13}}$.

 --Solution by \href{https://artofproblemsolving.comhttp://www.artofproblemsolving.com/Forum/memberlist.php?mode=viewprofile&u=200685}{TheMaskedMagician}

 


	\item (\textit{AHSME 1994}) A large rectangle is partitioned into four rectangles by two segments parallel to its sides. The areas of three of the resulting rectangles are shown. What is the area of the fourth rectangle?

                    \begin{figure}[H]
                        \centering
                        \includegraphics[width=0.5\textwidth]{Problem/AHSME_I_1994_2_Problem_Figure_1.png}
                    \end{figure}
                    
$\textbf{(A)}\ 10 \qquad\textbf{(B)}\ 15 \qquad\textbf{(C)}\ 20 \qquad\textbf{(D)}\ 21 \qquad\textbf{(E)}\ 25$

 \textbf{Answer: }B



\textbf{Solution}


                    \begin{figure}[H]
                        \centering
                        \includegraphics[width=0.5\textwidth]{Solution/AHSME_I_1994_2_Solution_Figure_1.png}
                    \end{figure}
                    

 We can easily see the dimensions of each small rectangle. So the area of the last rectangle is $3\times 5=\boxed{\textbf{(B) }15}$.

 --Solution by \href{https://artofproblemsolving.comhttp://www.artofproblemsolving.com/Forum/memberlist.php?mode=viewprofile&u=200685}{TheMaskedMagician}

 



\textbf{Solution 2}


                    \begin{figure}[H]
                        \centering
                        \includegraphics[width=0.5\textwidth]{Solution/AHSME_I_1994_2_Solution_Figure_2.png}
                    \end{figure}
                    

 In case its not immediately obvious from inspection what the dimensions of the small rectangles are, we can work it out.  We know $ax$ and $bx$ and $by$ and we want to know $ay$.   We can compute it as follows $ay = \frac{(ax)(by)}{bx} = \frac{ 6\cdot 35 }{14} = 15$ and the answer is $\fbox{B}$



	\item (\textit{AHSME 1994}) How many of the following are equal to $x^x+x^x$ for all $x>0$?


 $\textbf{I:}\ 2x^x \qquad\textbf{II:}\ x^{2x} \qquad\textbf{III:}\ (2x)^x \qquad\textbf{IV:}\ (2x)^{2x}$
$\textbf{(A)}\ 0 \qquad\textbf{(B)}\ 1 \qquad\textbf{(C)}\ 2 \qquad\textbf{(D)}\ 3 \qquad\textbf{(E)}\ 4$

 \textbf{Answer: }B



\textbf{Solution}

We look at each statement individually.

 $\textbf{I:}\ 2x^x$. We note that $x^x+x^x=x^x(1+1)=2x^x$. So statement $\textbf{I}$ is true.

 $\textbf{II:}\ x^{2x}$. We find a counter example which is $x=1$. $2\neq 1$. So statement $\textbf{II}$ is false.

 $\textbf{III:}\ (2x)^x$. We see that this statement is equal to $2^xx^x$. $x=2$ is a counter example. $8\neq 16$. So statement $\textbf{III}$ is false.

 $\textbf{IV:}\ (2x)^{2x}$. We see that $x=1$ is again a counter example. $2\neq 4$. So statement $\textbf{IV}$ is false.

 Therefore, our answer is $\boxed{\textbf{(B) }1}$.

 --Solution by \href{https://artofproblemsolving.comhttp://www.artofproblemsolving.com/Forum/memberlist.php?mode=viewprofile&u=200685}{TheMaskedMagician}

 


	\item (\textit{AHSME 1994}) In the $xy$-plane, the segment with endpoints $(-5,0)$ and $(25,0)$ is the diameter of a circle. If the point $(x,15)$ is on the circle, then $x=$
$\textbf{(A)}\ 10 \qquad\textbf{(B)}\ 12.5 \qquad\textbf{(C)}\ 15\qquad\textbf{(D)}\ 17.5 \qquad\textbf{(E)}\ 20$

 \textbf{Answer: }A



\textbf{Solution}

We see that the center of this circle is at $\left(\frac{-5+25}{2},0\right)=(10,0)$. The radius is $\frac{30}{2}=15$. So the equation of this circle is 
 \[(x-10)^2+y^2=225.\] Substituting $y=15$ yields $(x-10)^2=0$ so $x=\boxed{\textbf{(A) }10}$.

 --Solution by \href{https://artofproblemsolving.comhttp://www.artofproblemsolving.com/Forum/memberlist.php?mode=viewprofile&u=200685}{TheMaskedMagician}

 



\textbf{Solution 2}

The diameter of the circle is $25- (-5)=30$ so the radius is $15$ and the center of the circle is on the $x$-axis at $x=10$.  The only point with $y=15$ must be the point exactly above the center.

 


	\item (\textit{AHSME 1994}) Pat intended to multiply a number by $6$ but instead divided by $6$. Pat then meant to add $14$ but instead subtracted $14$. After these mistakes, the result was $16$. If the correct operations had been used, the value produced would have been


 $\textbf{(A)}\ \text{less than 400} \qquad\textbf{(B)}\ \text{between 400 and 600} \qquad\textbf{(C)}\ \text{between 600 and 800} \\ \textbf{(D)}\ \text{between 800 and 1000} \qquad\textbf{(E)}\ \text{greater than 1000}$

 \textbf{Answer: }E



\textbf{Solution}

We reverse the operations that he did and then use the correct operations. His end result is $16$. Before that, he subtracted $14$ which means that his number after the first operation was $30$. He divided by $6$ so his number was $180$. 

 Now, we multiply $180$ by $6$ to get $1080$. Finally, $1080+14=1094$. Since $1094>1000$, our answer is $\boxed{\textbf{(E)}\ \text{greater than 1000}}$.

 --Solution by \href{https://artofproblemsolving.comhttp://www.artofproblemsolving.com/Forum/memberlist.php?mode=viewprofile&u=200685}{TheMaskedMagician}

 


	\item (\textit{AHSME 1994}) In the sequence

 \[..., a, b, c, d, 0, 1, 1, 2, 3, 5, 8,...\]
each term is the sum of the two terms to its left. Find $a$.


 $\textbf{(A)}\ -3 \qquad\textbf{(B)}\ -1 \qquad\textbf{(C)}\ 0 \qquad\textbf{(D)}\ 1 \qquad\textbf{(E)}\ 3$

 \textbf{Answer: }A



\textbf{Solution}

We work backwards to find $a$. 

 
 \[d+0=1\implies d=1\] 
 \[c+1=0\implies c=-1\]
 \[b+(-1)=1\implies b=2\]
 \[a+2=-1\implies a=\boxed{\textbf{(A)}-3.}\]

 --Solution by \href{https://artofproblemsolving.comhttp://www.artofproblemsolving.com/Forum/memberlist.php?mode=viewprofile&u=200685}{TheMaskedMagician}

 


	\item (\textit{AHSME 1994}) Squares $ABCD$ and $EFGH$ are congruent, $AB=10$, and $G$ is the center of square $ABCD$. The area of the region in the plane covered by these squares is

                    \begin{figure}[H]
                        \centering
                        \includegraphics[width=0.5\textwidth]{Problem/AHSME_I_1994_7_Problem_Figure_1.png}
                    \end{figure}
                    
$\textbf{(A)}\ 75 \qquad\textbf{(B)}\ 100 \qquad\textbf{(C)}\ 125 \qquad\textbf{(D)}\ 150 \qquad\textbf{(E)}\ 175$

 \textbf{Answer: }E



\textbf{Solution}

The area of the entire region in the plane is the area of the figure. However, we cannot simply add the two areas of the squares. We find the area of $\triangle ABG$ and subtract this from $200$, the total area of the two squares.

 
Since $G$ is the center of $ABCD$, $BG$ is half of the diagonal of the square. The diagonal of $ABCD$ is $10\sqrt{2}$ so $BG=5\sqrt{2}$. Since $EFGH$ is a square, $\angle G=90^\circ$. So $\triangle ABG$ is an isosceles right triangle. Its area is $\frac{(5\sqrt{2})^2}{2}=\frac{50}{2}=25$. Therefore, the area of the region is $200-25=\boxed{\textbf{(E) }175.}$

 --Solution by \href{https://artofproblemsolving.comhttp://www.artofproblemsolving.com/Forum/memberlist.php?mode=viewprofile&u=200685}{TheMaskedMagician}

 



\textbf{Solution 2}

`Note the overlap area $\triangle ABG$ is exactly $1/4$ the area of $ABCD$.  So the total area is $100 + 100 - 100/4 = 175$



	\item (\textit{AHSME 1994}) In the polygon shown, each side is perpendicular to its adjacent sides, and all 28 of the sides are congruent. The perimeter of the polygon is $56$. The area of the region bounded by the polygon is

                    \begin{figure}[H]
                        \centering
                        \includegraphics[width=0.5\textwidth]{Problem/AHSME_I_1994_8_Problem_Figure_1.png}
                    \end{figure}
                    
$\textbf{(A)}\ 84 \qquad\textbf{(B)}\ 96 \qquad\textbf{(C)}\ 100 \qquad\textbf{(D)}\ 112 \qquad\textbf{(E)}\ 196$

 \textbf{Answer: }C



\textbf{Solution}

Since the perimeter is $56$ and all of the sides are congruent, the length of each side is $2$. We break the figure into squares as shown below.

 
                    \begin{figure}[H]
                        \centering
                        \includegraphics[width=0.5\textwidth]{Solution/AHSME_I_1994_8_Solution_Figure_1.png}
                    \end{figure}
                    

 We see that there are a total of $2(1+3+5)+7=25$ squares with side length $2$. Therefore, the total area is $4\cdot 25=\boxed{\textbf{(C) }100.}$

 --Solution by \href{https://artofproblemsolving.comhttp://www.artofproblemsolving.com/Forum/memberlist.php?mode=viewprofile&u=200685}{TheMaskedMagician}

 


	\item (\textit{AHSME 1994}) If $\angle A$ is four times $\angle B$, and the complement of $\angle B$ is four times the complement of $\angle A$, then $\angle B=$
$\textbf{(A)}\ 10^{\circ} \qquad\textbf{(B)}\ 12^{\circ} \qquad\textbf{(C)}\ 15^{\circ} \qquad\textbf{(D)}\ 18^{\circ} \qquad\textbf{(E)}\ 22.5^{\circ}$

 \textbf{Answer: }D



\textbf{Solution}

Let $\angle A=x$ and $\angle B=y$. From the first condition, we have $x=4y$. From the second condition, we have 
 \[90-y=4(90-x).\] Substituting $x=4y$ into the previous equation and solving yields 
 \begin{align*}90-y=4(90-4y)&\implies 90-y=360-16y\\&\implies 15y=270\\&\implies y=\boxed{\textbf{(D) }18^\circ.}\end{align*}


 --Solution by \href{https://artofproblemsolving.comhttp://www.artofproblemsolving.com/Forum/memberlist.php?mode=viewprofile&u=200685}{TheMaskedMagician}

 


	\item (\textit{AHSME 1994}) For distinct real numbers $x$ and $y$, let $M(x,y)$ be the larger of $x$ and $y$ and let $m(x,y)$ be the smaller of $x$ and $y$. If $a<b<c<d<e$, then

 \[M(M(a,m(b,c)),m(d,m(a,e)))=\]
$\textbf{(A)}\ a \qquad\textbf{(B)}\ b \qquad\textbf{(C)}\ c \qquad\textbf{(D)}\ d \qquad\textbf{(E)}\ e$

 \textbf{Answer: }B



\textbf{Solution}

We work thorough the equation step by step, simplifying as follows:

 
 \begin{align*}M(M(a,m(b,c)),m(d,m(a,e)))&=M(M(a,b),m(d,a))\\&=M(b,a)\\&=\boxed{\textbf{(B) }b}\end{align*}


 --Solution by \href{https://artofproblemsolving.comhttp://www.artofproblemsolving.com/Forum/memberlist.php?mode=viewprofile&u=200685}{TheMaskedMagician}

 


	\item (\textit{AHSME 1994}) Three cubes of volume $1, 8$ and $27$ are glued together at their faces. The smallest possible surface area of the resulting configuration is


 $\textbf{(A)}\ 36 \qquad\textbf{(B)}\ 56 \qquad\textbf{(C)}\ 70 \qquad\textbf{(D)}\ 72 \qquad\textbf{(E)}\ 74$

 \textbf{Answer: }D


	\item (\textit{AHSME 1994}) If $i^2=-1$, then $(i-i^{-1})^{-1}=$
$\textbf{(A)}\ 0 \qquad\textbf{(B)}\ -2i \qquad\textbf{(C)}\ 2i \qquad\textbf{(D)}\ -\frac{i}{2} \qquad\textbf{(E)}\ \frac{i}{2}$

 \textbf{Answer: }D



\textbf{Solution}

We simplify step by step as follows: 
 \begin{align*}(i-i^{-1})^{-1}&=\frac{1}{i-i^{-1}}\\&=\frac{1}{i-\frac{1}{i}}\\&=\frac{1}{\left(\frac{i^2-1}{i}\right)}\\&=\frac{i}{i^2-1}\\&=\boxed{\textbf{(D) }-\frac{i}{2}.}\end{align*}


 --Solution by \href{https://artofproblemsolving.comhttp://www.artofproblemsolving.com/Forum/memberlist.php?mode=viewprofile&u=200685}{TheMaskedMagician}

 


	\item (\textit{AHSME 1994}) In triangle $ABC$, $AB=AC$. If there is a point $P$ strictly between $A$ and $B$ such that $AP=PC=CB$, then $\angle A =$

                    \begin{figure}[H]
                        \centering
                        \includegraphics[width=0.5\textwidth]{Problem/AHSME_I_1994_13_Problem_Figure_1.png}
                    \end{figure}
                    
$\textbf{(A)}\ 30^{\circ} \qquad\textbf{(B)}\ 36^{\circ} \qquad\textbf{(C)}\ 48^{\circ} \qquad\textbf{(D)}\ 60^{\circ} \qquad\textbf{(E)}\ 72^{\circ}$

 \textbf{Answer: }B



\textbf{Solution}


                    \begin{figure}[H]
                        \centering
                        \includegraphics[width=0.5\textwidth]{Solution/AHSME_I_1994_13_Solution_Figure_1.png}
                    \end{figure}
                    

 Let $\angle A=x$. Since $AP=PC$, we have $\angle ACP=x$ as well. Then $\angle APC=180-2x\implies\angle BPC=\angle CBP=2x$. Since $AB=AC$, we have $\angle CBP=\frac{180-x}{2}$. 

 So $2x=\frac{180-x}{2}\implies 5x=180\implies x=\angle A=\boxed{\textbf{(B) }36^\circ.}$

 
--Solution by \href{https://artofproblemsolving.comhttp://www.artofproblemsolving.com/Forum/memberlist.php?mode=viewprofile&u=200685}{TheMaskedMagician}

 


	\item (\textit{AHSME 1994}) Find the sum of the arithmetic series

 \[20+20\frac{1}{5}+20\frac{2}{5}+\cdots+40\]
$\textbf{(A)}\ 3000 \qquad\textbf{(B)}\ 3030 \qquad\textbf{(C)}\ 3150 \qquad\textbf{(D)}\ 4100 \qquad\textbf{(E)}\ 6000$

 \textbf{Answer: }B



\textbf{Solution}

\textbf{Brief Introduction} For those that do not know the formula, the sum of an arithmetic series with first term $a_1$, last term $a_n$ as $n$ terms, is 
 \[S = \frac{n(a_1+a_n)}{2}.\] We can prove this as follows:

 Let $d$ be the common difference between terms of our series and let $n$ be the number of terms in our series. Let $a_1$ be the first term. Our series is 
 \[a_1,~a_1+d,~a_1+2d,\dots,a_1+(n-1)d.\] Note that we have $n-1$ in the last term because $a_1$ is a term. Let $S$ be our sum such that 
 \[S=a_1+(a_1+d)+\dots+(a_1+(n-1)d).\] We can rewrite our sums as 
 \[S=(a_1+(n-1)d)+(a_1+(n-2)d)+\dots+(a_1+d)+a_1.\] Adding these two sums together essentially creates $n$ pairs of $a_1+(a_1+(n-1)d)$ as shown below: 
 \[2S=n(a_1+(a_1+(n-1)d))\implies S=\frac{n(a_1+a_n)}{2}\] We use $a_n$ in place of $a_1+(n-1)d$ to represent the last term.

 


 \textbf{Solving} Our first term is $20$ and our last term is $40$. To find the number of terms, $n$, we note that the common difference between each term is $\frac{1}{5}$. So we have 
 \[20+\frac{1}{5}(n-1)=40\implies n-1=100\implies n=101.\] Using our formula, our sum is 
 \[101\left(\frac{20+40}{2}\right)=101\times 30=\boxed{\textbf{(B) }3030.}\]

 
--Solution by \href{https://artofproblemsolving.comhttp://www.artofproblemsolving.com/Forum/memberlist.php?mode=viewprofile&u=200685}{TheMaskedMagician}

 


	\item (\textit{AHSME 1994}) For how many $n$ in $\{1, 2, 3, ..., 100 \}$ is the tens digit of $n^2$ odd?


 $\textbf{(A)}\ 10 \qquad\textbf{(B)}\ 20 \qquad\textbf{(C)}\ 30 \qquad\textbf{(D)}\ 40 \qquad\textbf{(E)}\ 50$

 \textbf{Answer: }B



\textbf{Solution}

Let $n=10a+b$. So $n^2=(10a+b)^2=100a^2+20ab+b^2$. The term $100a^2$ only contributes digits starting at the hundreds place, so this does not affect whether the tens digit is odd.  The term $20ab$ only contributes digits starting at the tens place, and the tens digit contributed will be the ones digit of $2ab$ which is even.  So we see this term also does not affect whether the tens digit is odd.  This means only $b^2$ can affect whether the tens digit is odd.  We can quickly check $1^2=1, \dots, 9^2=81$ and discover only for $b=4$ or $b=6$ does $b^2$ have an odd tens digit. The total number of positive integers less than or equal to $100$ that have $4$ or $6$ as the units digit is 
 \[10\times 2=\boxed{\textbf{(B) }20.}\]

 --Solution by \href{https://artofproblemsolving.comhttp://www.artofproblemsolving.com/Forum/memberlist.php?mode=viewprofile&u=200685}{TheMaskedMagician}

 


	\item (\textit{AHSME 1994}) Some marbles in a bag are red and the rest are blue. If one red marble is removed, then one-seventh of the remaining marbles are red. If two blue marbles are removed instead of one red, then one-fifth of the remaining marbles are red. How many marbles were in the bag originally?


 $\textbf{(A)}\ 8 \qquad\textbf{(B)}\ 22 \qquad\textbf{(C)}\ 36 \qquad\textbf{(D)}\ 57 \qquad\textbf{(E)}\ 71$

 \textbf{Answer: }B



\textbf{Solution}

Let $r$ and $b$ be the number of red and blue marbles originally in the bag respectively. After $1$ red marble is removed, there are $r+b-1$ marbles left in the bag and $r-1$ red marbles left. So 
 \[\frac{r-1}{r+b-1}=\frac{1}{7}.\] When $2$ blue marbles are removed, there are $r$ red marbles and $r+b-2$ total marbles left in the bag. So 
 \[\frac{r}{r+b-2}=\frac{1}{5}.\] Cross multiplying for each yields 
 \begin{align*}7r-7=r+b-1&\implies 7r-6=r+b\\ 5r=r+b-2&\implies 5r+2=r+b.\end{align*}
 We can equate each of these expressions to yields 
 \[7r-6=5r+2\implies 2r=8\implies r=4\implies b=18.\] Therefore, the total number of marbles is 
 \[r+b=4+18=\boxed{\textbf{(B) }22.}\]

 --Solution by \href{https://artofproblemsolving.comhttp://www.artofproblemsolving.com/Forum/memberlist.php?mode=viewprofile&u=200685}{TheMaskedMagician}

 


	\item (\textit{AHSME 1994}) An $8$ by $2\sqrt{2}$ rectangle has the same center as a circle of radius $2$. The area of the region common to both the rectangle and the circle is


 $\textbf{(A)}\ 2\pi \qquad\textbf{(B)}\ 2\pi+2 \qquad\textbf{(C)}\ 4\pi-4 \qquad\textbf{(D)}\ 2\pi+4 \qquad\textbf{(E)}\ 4\pi-2$

 \textbf{Answer: }D



\textbf{Solution}


                    \begin{figure}[H]
                        \centering
                        \includegraphics[width=0.5\textwidth]{Solution/AHSME_I_1994_17_Solution_Figure_1.png}
                    \end{figure}
                    

 We draw the diagram above. Dropping an altitude from the center of the circle to the top side of the rectangle, yields a segment of length $\sqrt{2}$. Since the hypotenuse is the radius of the circle, it has length 2 and by Pythagoras the other leg of the triangle also has length $\sqrt{2}$.  We deduce that the two triangles formed by two radii of circle $O$ and the segment of the rectangle are 45-45-90 triangles. 

 The overlapping area consists of two sectors with central angle 90 and four 45-45-90 triangles with base $\sqrt{2}$ and height $\sqrt{2}$.  The area of the circle is $4\pi$ so the area of the sectors is $2 \cdot 4\pi/4 = 2\pi$.  The area of the triangles is  $4\cdot (\sqrt{2})^2 /2 = 4$.  The combined area is 
 \[\boxed{\textbf{(D) }2\pi+4.}\]

 --Solution by \href{https://artofproblemsolving.comhttp://www.artofproblemsolving.com/Forum/memberlist.php?mode=viewprofile&u=200685}{TheMaskedMagician}

 


	\item (\textit{AHSME 1994}) Triangle $ABC$ is inscribed in a circle, and $\angle B = \angle C = 4\angle A$. If $B$ and $C$ are adjacent vertices of a regular polygon of $n$ sides inscribed in this circle, then $n=$

                    \begin{figure}[H]
                        \centering
                        \includegraphics[width=0.5\textwidth]{Problem/AHSME_I_1994_18_Problem_Figure_1.png}
                    \end{figure}
                    
$\textbf{(A)}\ 5 \qquad\textbf{(B)}\ 7 \qquad\textbf{(C)}\ 9 \qquad\textbf{(D)}\ 15 \qquad\textbf{(E)}\ 18$

 \textbf{Answer: }C



\textbf{Solution}

We solve for $\angle A$ as follows: 
 \[4\angle A+4\angle A+\angle A=180\implies 9\angle A=180\implies \angle A=20^\circ.\] That means that minor arc $\widehat{BC}$ has measure $40^\circ$. We can fit a maximum of $\frac{360}{40}=\boxed{\textbf{(C) }9}$ of these arcs in the circle.

 --Solution by \href{https://artofproblemsolving.comhttp://www.artofproblemsolving.com/Forum/memberlist.php?mode=viewprofile&u=200685}{TheMaskedMagician}

 


	\item (\textit{AHSME 1994}) Label one disk "$1$", two disks "$2$", three disks "$3$"$, ...,$ fifty disks "$50$". Put these $1+2+3+ \cdots+50=1275$ labeled disks in a box. Disks are then drawn from the box at random without replacement. The minimum number of disks that must be drawn to guarantee drawing at least ten disks with the same label is


 $\textbf{(A)}\ 10 \qquad\textbf{(B)}\ 51 \qquad\textbf{(C)}\ 415 \qquad\textbf{(D)}\ 451 \qquad\textbf{(E)}\ 501$

 \textbf{Answer: }C



\textbf{Solution}

We can solve this problem by thinking of the worst case scenario, essentially an adaptation of the Pigeon-hole principle. We can start by picking up all the disks numbered 1 to 9 since even if we have all those disks we won't have 10 of any one disk. This gives us 45 disks. 

 From disks numbered from 10 to 50, we can pick up at most 9 disks to prevent picking up 10. There are 50-10+1 = 41 different numbers from 10 to 50. We pick up 9 from each number, therefore, we multiply $41 \cdot 9 = 369$. In total, the maximum number we can pick up without picking up 10 of the same kind is $369+45=414$. We need one more disk to guarantee a complete set of 10. Therefore, the answer is $\boxed{415}$.

 


	\item (\textit{AHSME 1994}) Suppose $x,y,z$ is a geometric sequence with common ratio $r$ and $x \neq y$. If $x, 2y, 3z$ is an arithmetic sequence, then $r$ is


 $\textbf{(A)}\ \frac{1}{4} \qquad\textbf{(B)}\ \frac{1}{3} \qquad\textbf{(C)}\ \frac{1}{2} \qquad\textbf{(D)}\ 2 \qquad\textbf{(E)}\ 4$

 \textbf{Answer: }B



\textbf{Solution}

Let $y=xr, z=xr^2$. Since $x, 2y, 3z$ are an arithmetic sequence, there is a common difference and we have $2xr-x=3xr^2-2xr$. Dividing through by $x$, we get $2r-1=3r^2-2r$ or, rearranging, $(r-1)(3r-1)=0$. Since we are given $x\neq y\implies r\neq 1$, the answer is $\boxed{\textbf{(B)}\ \frac{1}{3}}$



	\item (\textit{AHSME 1994}) Find the number of counter examples to the statement:

 \[``\text{If  N is an odd positive integer the sum of whose digits is 4 and none of whose digits is 0, then N is prime}."\]
$\textbf{(A)}\ 0 \qquad\textbf{(B)}\ 1 \qquad\textbf{(C)}\ 2 \qquad\textbf{(D)}\ 3 \qquad\textbf{(E)}\ 4$

 \textbf{Answer: }C



\textbf{Solution}

Since the sum of the digits of $N$ is $4$ and none of the digits are $0$, $N$'s digits must the elements of one of the sets $\{1,1,1,1\},\{1,1,2\}$, $\{2,2\}$, $\{1,3\}$, or $\{4\}$. 

 In the first case, $N = 1111 = 101 \cdot 11$ so this is a counter example.  

 In the second case, $N=112$ is excluded for being even.  With $N=121=11^2$ we have a counterexample.  We can check $N=211$ by trial division, and verify it is indeed prime. 

 In the third case, $N=22$ is excluded for being even.

 In the fourth case, both $N=13$ and $N=31$ are prime.  

 In the last case $N=4$ is excluded for being even.

 This gives two counterexamples and the answer is $\fbox{C}$



	\item (\textit{AHSME 1994}) Nine chairs in a row are to be occupied by six students and Professors Alpha, Beta and Gamma. These three professors arrive before the six students and decide to choose their chairs so that each professor will be between two students. In how many ways can Professors Alpha, Beta and Gamma choose their chairs?


 $\textbf{(A)}\ 12 \qquad\textbf{(B)}\ 36 \qquad\textbf{(C)}\ 60 \qquad\textbf{(D)}\ 84 \qquad\textbf{(E)}\ 630$

 \textbf{Answer: }C



\textbf{Solution}

Since each professor must sit between two students, they cannot be seated in seats $1$ or $9$ (the seats at either end of the row). Hence, each professor has $7$ seats they can choose from and must be at least $1$ seat apart.

 This question is equivalent to choosing $3$ seats from a row of $5$ seats with no restrictions because we can simply generate a valid arrangement by inserting a seat right after the first and second seat chosen.

 Hence, the answer is $\binom{5}{3} \cdot 3! =$ $60$ $\textbf{(C)}$.

 


	\item (\textit{AHSME 1994}) In the $xy$-plane, consider the L-shaped region bounded by horizontal and vertical segments with vertices at $(0,0), (0,3), (3,3), (3,1), (5,1)$ and $(5,0)$. The slope of the line through the origin that divides the area of this region exactly in half is

                    \begin{figure}[H]
                        \centering
                        \includegraphics[width=0.5\textwidth]{Problem/AHSME_I_1994_23_Problem_Figure_1.png}
                    \end{figure}
                    
$\textbf{(A)}\ \frac{2}{7} \qquad\textbf{(B)}\ \frac{1}{3} \qquad\textbf{(C)}\ \frac{2}{3} \qquad\textbf{(D)}\ \frac{3}{4} \qquad\textbf{(E)}\ \frac{7}{9}$

 \textbf{Answer: }E



\textbf{Solution}

Let the vertices be $A=(0,0),B=(0,3),C=(3,3),D=(3,1),E=(5,1),F=(5,0)$. It is easy to see that the line must pass through $CD$. Let the line intersect $CD$ at the point $G=(3,3-x)$ (i.e. the point $x$ units below $C$). Since the quadrilateral $ABCG$ and pentagon $GDEFA$ must have the same area, we have the equation $3\times\frac{1}{2}\times(x+3)=\frac{1}{2}\times3\times(3-x)+2$. This simplifies into $3x=2$, or $x=\frac{2}{3}$, so $G=(3,\frac{7}{3})$. Therefore the slope of the line is $\boxed{\textbf{(E)}\ \frac{7}{9}}$




\textbf{Solution 2}

Consider the small rectangle between $x=3$ and $x=5$ with area $2$.  If we exclude that area then the shape becomes a much simpler, its just a rectangle.  In order to offset the exclusion of that area from the shape below the line, we must also exclude a shape with the same area above the line.  We want the remainder after excluding both areas to be  simple, so exclude a rectangle between $x=0$ and $x=3$ and $y=3$ and $y=c$ for some unknown $c$.  For the area to be the same we need $(3-c) \cdot 3 = 2$, or $c=7/3$.  After excluding our two offsetting areas, we're left with a rectangle from $x=0$ to $x=3$ and $y=0$ to $y=c$.  The area of this region is clearly bisected by its diagonal line.  The line passes through $(0,0)$ and $(3,c)$ so its slope is $c/3 = 7/9$ and the answer is $\fbox{E}$



	\item (\textit{AHSME 1994}) A sample consisting of five observations has an arithmetic mean of $10$ and a median of $12$. The smallest value that the range (largest observation minus smallest) can assume for such a sample is


 $\textbf{(A)}\ 2 \qquad\textbf{(B)}\ 3 \qquad\textbf{(C)}\ 5 \qquad\textbf{(D)}\ 7 \qquad\textbf{(E)}\ 10$

 \textbf{Answer: }C



\textbf{Solution}

The minimum range occurs in the set $\{7,7,12,12,12\}$, so the answer is $\boxed{\textbf{(C)}\ 5}$



	\item (\textit{AHSME 1994}) If $x$ and $y$ are non-zero real numbers such that

 \[|x|+y=3 \qquad \text{and} \qquad |x|y+x^3=0,\]
then the integer nearest to $x-y$ is


 $\textbf{(A)}\ -3 \qquad\textbf{(B)}\ -1 \qquad\textbf{(C)}\ 2 \qquad\textbf{(D)}\ 3 \qquad\textbf{(E)}\ 5$

 \textbf{Answer: }A



\textbf{Solution}

We have two cases to consider: $x$ is positive or $x$ is negative. If $x$ is positive, we have $x+y=3$ and $xy+x^3=0$

 Solving for $y$ in the top equation gives us $3-x$. Plugging this in gives us: $x^3-x^2+3x=0$.  Since we're told $x$ is not zero, we can divide by $x$, giving us: $x^2-x+3=0$

 The discriminant of this is $(-1)^2-4(1)(3)=-11$, which means the equation has no real solutions. 

 We conclude that $x$ is negative. In this case $-x+y=3$ and $-xy+x^3=0$.  Negating the top equation gives us $x-y=-3$. We seek $x-y$, so the answer is $\boxed{(A) -3}$

 -solution by jmania

 


	\item (\textit{AHSME 1994}) A regular polygon of $m$ sides is exactly enclosed (no overlaps, no gaps) by $m$ regular polygons of $n$ sides each. (Shown here for $m=4, n=8$.) If $m=10$, what is the value of $n$?

                    \begin{figure}[H]
                        \centering
                        \includegraphics[width=0.5\textwidth]{Problem/AHSME_I_1994_26_Problem_Figure_1.png}
                    \end{figure}
                    
$\textbf{(A)}\ 5 \qquad\textbf{(B)}\ 6 \qquad\textbf{(C)}\ 14 \qquad\textbf{(D)}\ 20 \qquad\textbf{(E)}\ 26$

 \textbf{Answer: }A



\textbf{Solution}

To find the number of sides on the regular polygons that surround the decagon, we can find the interior angles and work from there. Knowing that the measure of the interior angle of any regular polygon is $\frac{(n-2)*180}{n}$, the measure of the decagon's interior angle is $\frac{8*180}{10} = 144$ degrees. 

 
The regular polygons meet at every vertex such that the angle outside of the decagon is divided evenly in two. With this, we know that the angle of the regular polygon is $216/2=108$ degrees. Using the previous formula, $n=5$ $\boxed{\textbf{(A) }5}$



	\item (\textit{AHSME 1994}) A bag of popping corn contains $\frac{2}{3}$ white kernels and $\frac{1}{3}$ yellow kernels. Only $\frac{1}{2}$ of the white kernels will pop, whereas $\frac{2}{3}$ of the yellow ones will pop. A kernel is selected at random from the bag, and pops when placed in the popper. What is the probability that the kernel selected was white?


 $\textbf{(A)}\ \frac{1}{2} \qquad\textbf{(B)}\ \frac{5}{9} \qquad\textbf{(C)}\ \frac{4}{7} \qquad\textbf{(D)}\ \frac{3}{5} \qquad\textbf{(E)}\ \frac{2}{3}$

 \textbf{Answer: }D



\textbf{Solution}

To find the probability that the kernel is white, the probability of $P(\mathrm{white|popped}) = \frac{P(\mathrm{white, popped})}{P(\mathrm{popped})}$.

 Running a bit of calculations $P(\mathrm{white, popped}) = \frac{1}{3}$ while $P(\mathrm{popped}) = \frac{1}{3} + \frac{2}{9} = \frac{5}{9}$. Plugging this into the earlier equation, $P(\mathrm{white|popped}) = \frac{\frac{1}{3}}{\frac{5}{9}}$, meaning that the answer is $\boxed{\textbf{(D)}\ \frac{3}{5}}$.

 


	\item (\textit{AHSME 1994}) In the $xy$-plane, how many lines whose $x$-intercept is a positive prime number and whose $y$-intercept is a positive integer pass through the point $(4,3)$?


 $\textbf{(A)}\ 0 \qquad\textbf{(B)}\ 1 \qquad\textbf{(C)}\ 2 \qquad\textbf{(D)}\ 3 \qquad\textbf{(E)}\ 4$

 \textbf{Answer: }C



\textbf{Solution 1}

The line with $x$-intercept  $a$ and $y$-intercept $b$ is given by the equation $\frac{x}{a} + \frac{y}{b} = 1$.  We are told $(4,3)$ is on the line so 

 
 \[\frac{4}{a} + \frac{3}{b} = 1 \implies ab - 4b - 3a = 0 \implies (a-4)(b-3)=12\]

 Since $a$ and $b$ are integers, this equation holds only if $(a-4)$ is a factor of $12$.  The factors are $1, 2, 3, 4, 6, 12$ which means $a$ must be one of $5, 6, 7, 8, 10, 16$.  The only members of this list which are prime are $a=5$ and $a=7$, so the number of solutions is $\boxed{\textbf{(C) } 2}$.

 



\textbf{Solution 2}


                    \begin{figure}[H]
                        \centering
                        \includegraphics[width=0.5\textwidth]{Solution/AHSME_I_1994_28_Solution_Figure_1.png}
                    \end{figure}
                    Let $C = (4,3)$, $DF=a$, and $AD=b$. As stated in the problem, the $x$-intercept $DF=a$ is a positive prime number, and the $y$-intercept $AD=b$ is a positive integer.

 Through similar triangles, $\frac{AB}{BC}=\frac{CE}{EF}$, $\frac{b-3}{4}=\frac{3}{a-4}$, $(a-4)(b-3)=12$

 The only cases where $a$ is prime are:
 \[\begin{cases}  a-4=1 & a=5 \\ b-3=12 & b=15 \end{cases}\]

 \[and\]

 \[\begin{cases}  a-4=3 & a=7 \\ b-3=4 & b=5 \end{cases}\]

 So the number of solutions are $\boxed{\textbf{(C) }2}$.

 


	\item (\textit{AHSME 1994}) Points $A, B$ and $C$ on a circle of radius $r$ are situated so that $AB=AC, AB>r$, and the length of minor arc $BC$ is $r$. If angles are measured in radians, then $AB/BC=$

                    \begin{figure}[H]
                        \centering
                        \includegraphics[width=0.5\textwidth]{Problem/AHSME_I_1994_29_Problem_Figure_1.png}
                    \end{figure}
                    
$\textbf{(A)}\ \frac{1}{2}\csc{\frac{1}{4}} \qquad\textbf{(B)}\ 2\cos{\frac{1}{2}} \qquad\textbf{(C)}\ 4\sin{\frac{1}{2}} \qquad\textbf{(D)}\ \csc{\frac{1}{2}} \qquad\textbf{(E)}\ 2\sec{\frac{1}{2}}$

 \textbf{Answer: }A



\textbf{Solution 1}

First note that arc length equals $r\theta$, where $\theta$ is the central angle in radians. Call the center of the circle $O$. Then $\angle{BOC} = 1$ radian because the minor arc $BC$ has length $r$. Since $ABC$ is isosceles, $\angle{AOB} = \pi - \frac{1}{2}$. We use the Law of Cosines to find that 
 \[\frac{AB}{BC} = \frac{\sqrt{2r^2 - 2r^2\cos{(\pi - \frac{1}{2})}}}{\sqrt{2r^2 - 2r^2\cos1}} = \frac{\sqrt{1 + \cos{(\frac{1}{2})}}}{\sqrt{1 - \cos1}}.\] Using half-angle formulas, we have that this ratio simplifies to 
 \[\frac{\cos\frac{1}{4}}{\sin{\frac{1}{2}}} = \frac{\cos\frac{1}{4}}{\sqrt{1 - \cos^2{\frac{1}{2}}}} = \frac{\cos\frac{1}{4}}{\sqrt{(1 + \cos{\frac{1}{2}})(1 - \cos{\frac{1}{2}})}} = \frac{\cos{\frac{1}{4}}}{2\cos{\frac{1}{4}}\sin{\frac{1}{4}}}\] 
 \[= \boxed{\frac{1}{2}\csc{\frac{1}{4}}.}\]




\textbf{Solution 2}


                    \begin{figure}[H]
                        \centering
                        \includegraphics[width=0.5\textwidth]{Solution/AHSME_I_1994_29_Solution_Figure_1.png}
                    \end{figure}
                    

 Let the center of this circle be $O$, $\angle BOC = \theta$, the radius of $\odot O$ be $r$.

 By the definition of radian, $\theta =$$\overarc{BC}$$/r=1$
$\angle BAC=\frac{\angle BOC}{2} = \frac12$
$\sin \angle BAD = \frac{BD}{AB}$, $\sin \frac{\angle BAC}{2} = \frac{\frac{BC}{2}}{AB}$
$\frac{BC}{AB}=2 \sin \frac{\frac12}{2} = 2\sin \frac14$
$\frac{AB}{BC}= \frac{1}{2\sin \frac14} = \boxed{\textbf{(A) } \frac{1}{2} \csc \frac{1}{4} }$



	\item (\textit{AHSME 1994}) When $n$ standard 6-sided dice are rolled, the probability of obtaining a sum of 1994 is greater than zero and is the same as the probability of obtaining a sum of $S$. The smallest possible value of $S$ is


 $\textbf{(A)}\ 333 \qquad\textbf{(B)}\ 335 \qquad\textbf{(C)}\ 337 \qquad\textbf{(D)}\ 339 \qquad\textbf{(E)}\ 341$

 \textbf{Answer: }C



\textbf{Solution}

Let $d_i$ be the number on the $i$th die.  There is a symmetry where we can replace each die's number with $d_i' = 7-d_i$.  Note that applying the symmetry twice we get back to where we started since $7-(7-d_i)=d_i$, so this symmetry is its own inverse.  Under this symmetry the sum $R=\sum_{i=1}^n d_i$ is replaced by $S = \sum_{i=1}^n 7-d_i = 7n - R$.  As a result of this symmetry the sum $R$ the sum $S$ have the same probability because any combination of $d_i$ which sum to $R$ can be replaced with $d_i'$ which sum to $S$, and conversely.  In other words, there is a one-to-one mapping between the combinations of dice which sum to $R$ and the combinations which sum to $S$.

 When $R=1994$ we seek the smallest number $S=7n-1994$, which clearly happens when $n$ is smallest.  Therefore we want to find the smallest $n$ which gives non-zero probability of obtaining $R=1994$.  This occurs when there are just enough dice for this sum to be possible, and any fewer dice would result in $R=1994$ being impossible.  Thus $n = \left\lceil \frac{1994}{6} \right\rceil = 333$ and $S = 333\cdot 7 - 1994 = 337$.  The answer is $\textbf{(C)}$.

 


\end{enumerate}